\documentclass[usenames,dvipsnames]{beamer}

\usefonttheme{serif}
\setbeamertemplate{items}[circle]

\usepackage{hyperref}
\hypersetup{pdfborderstyle={/S/U/W 1}}
\usepackage{mathpazo}
\usepackage{pxfonts}
\usepackage{tcolorbox}
\usepackage{soul}
\usepackage{tikz}
\usepackage{algorithm}
\usepackage{algpseudocode}

\beamertemplatenavigationsymbolsempty

\DeclareMathOperator*{\Exp}{\mathbf{E}}
\DeclareMathOperator{\Regret}{Regret}
\DeclareMathOperator*{\argmin}{argmin}
\DeclareMathOperator*{\argmax}{argmax}
\DeclareMathOperator*{\maximize}{maximize}
\DeclareMathOperator*{\minimize}{minimize}
\DeclareMathOperator*{\subjectto}{subject\ to}

\newcommand{\R}{\mathbb{R}}
\newcommand{\indicator}{\mathbf{1}}
\newcommand{\norm}[1]{\left\|#1\right\|}

\newcommand{\Cite}[1]{{\tiny \textcolor{Blue}{[#1]}}}

\newtheorem{proposition}{Proposition}

\makeatletter
\newcommand\SoulColor{%
\let\set@color\beamerorig@set@color
\let\reset@color\beamerorig@reset@color}
\makeatother
\setstcolor{red}

% TikZ styles for matching market diagrams
\definecolor{babyblue}{RGB}{137,207,240}
\tikzset{
    matching/.style={
        every node/.style={circle, draw, thick, minimum size=0.7cm, inner sep=0pt},
        man/.style={fill=babyblue},
        woman/.style={fill=pink!60},
        edge/.style={ultra thick},
        matchedge/.style={ultra thick, red},
    },
}

% Helpers for the Gale-Shapley animation: highlight a woman's name in a man's
% preference list once she has been proposed to by him.
\newcommand{\prefw}[1]{\colorbox{white}{$#1$}}
\newcommand{\proposed}[2]{\alt<#1->{\colorbox{yellow!85}{$#2$}}{\prefw{#2}}}


\title{Pacing}

\date{August 2026}
\author{D\'avid P\'al}
\institute{Uber, Sunnyvale}

\begin{document}

%%%%%%%%%%%%%%%%%%%%%%%%%%%%%%%%%%%%%%%%%%%%%%%%%%%%%%%%%%%%%%%%%%%%%%%%%%%%%%%%%%%%%%%%%%
\begin{frame}
\maketitle
\end{frame}

%%%%%%%%%%%%%%%%%%%%%%%%%%%%%%%%%%%%%%%%%%%%%%%%%%%%%%%%%%%%%%%%%%%%%%%%%%%%%%%%%%%%%%%%%%
\begin{frame}
\frametitle{Online Decision Making}

\begin{itemize}
\item Requests $R_1, R_2, \dots, R_T$ arrive, one at a time.
\item For each request $R_t$, we need to choose an action $I_t \in \{1,2,\dots,N\}$.
\item Each action $i$, has a vector of predicted business metrics
$$
(R_{t,i,0}, R_{t,i,1}, \dots, R_{t,i,N}) \in \R^{M+1}
$$
associated with the request $R_t$.
\item In other words, the action $R_t$ is an $N \times (M+1)$ matrix.
\item The sequence of actions $I_1, I_2, \dots, I_T$ should maximize a
certain business metric subject to constraints on other business metrics.
\end{itemize}
\end{frame}

%%%%%%%%%%%%%%%%%%%%%%%%%%%%%%%%%%%%%%%%%%%%%%%%%%%%%%%%%%%%%%%%%%%%%%%%%%%%%%%%%%%%%%%%%%
\begin{frame}
\frametitle{Linear Program Formulation}

We encode action $I_t$ as a vector $x_t \in \R^N$

\begin{align*}
\maximize_{x_{1,1}, \dots, x_{T,N} \in \R} \quad & \frac{1}{T} \sum_{t=1}^T \sum_{i=1}^N R_{t,i,0} \cdot x_{t,i} \\
\subjectto \quad & \frac{1}{T} \sum_{t=1}^T \sum_{i=1}^N R_{t,i,j} \cdot x_{t,i} \le B_j && \quad 1 \le j \le M \\
& \sum_{i=1}^N x_{t,i} = 1 && \quad 1 \le t \le T \\
& x_{t,i} \ge 0 && \quad 1 \le i \le N, 1 \le t \le T
\end{align*}

We allow randomization.
\end{frame}

%%%%%%%%%%%%%%%%%%%%%%%%%%%%%%%%%%%%%%%%%%%%%%%%%%%%%%%%%%%%%%%%%%%%%%%%%%%%%%%%%%%%%%%%%%
\begin{frame}
\frametitle{Policy Optimization}

\begin{itemize}
\item Equivalently, the problem can be formulated as a policy optimization problem.
\item Define the probability simplex
$$
\Delta_N = \left\{ (x_1, x_2, \dots, x_N) \in \R^N ~:~ x_1, x_2, \dots, x_N \ge 0, \sum_{i=1}^N x_i = 1  \right\} \: .
$$
\item Find a policy $\pi: \R^{N \times (M+1)} \to \Delta_N$ such that
\end{itemize}

\begin{align*}
\maximize_{\pi \in \Pi} \quad & \frac{1}{T} \sum_{t=1}^T \sum_{i=1}^N R_{t,i,0} \cdot \pi(R_t)_i \\
\subjectto \quad & \frac{1}{T} \sum_{t=1}^T \sum_{i=1}^N R_{t,i,j} \cdot \pi(R_t)_i \le B_j & 1 \le j \le M
\end{align*}

\end{frame}

%%%%%%%%%%%%%%%%%%%%%%%%%%%%%%%%%%%%%%%%%%%%%%%%%%%%%%%%%%%%%%%%%%%%%%%%%%%%%%%%%%%%%%%%%%
\begin{frame}
\frametitle{Distributional Policy Optimization}

\begin{itemize}
\item We don't know the sequence of requests $R_1, R_2, \dots, R_T$ ahead of time.
\item So, we can't solve the linear program or the policy optimization problem directly.
\item We make a simple assumption that the requests are sampled i.i.d. from a fixed, but unknown, distribution $D$.
\item Find a policy $\pi: \R^{N \times (M+1)} \to \Delta_N$ such that
\end{itemize}

\begin{align*}
\maximize_{\pi \in \Pi} \quad & \Exp_{R \sim D} \left[ \sum_{i=1}^N \pi(R)_i \cdot R_{i,0} \right] \\
\subjectto \quad & \Exp_{R \sim D} \left[ \sum_{i=1}^N \pi(R)_i \cdot R_{i,j} \right] \le B_j & 1 \le j \le M
\end{align*}

\begin{itemize}
\item We assume that we have a historical i.i.d. sample from $D$ available.
\end{itemize}

\end{frame}

%%%%%%%%%%%%%%%%%%%%%%%%%%%%%%%%%%%%%%%%%%%%%%%%%%%%%%%%%%%%%%%%%%%%%%%%%%%%%%%%%%%%%%%%%%
\begin{frame}
\frametitle{Lagrangian Policies}

\begin{theorem}[Lagrangian Policy]
Let $D$ be a distribution on $\R^{N \times (M+1)}$ with a probability density.
The distributional policy optimization problem has an optimal solution of the
form $\pi_\lambda:\R^{N \times (M+1)} \to \Delta_N$ given by
$$
\pi_\lambda(R) = e_{i^*} \qquad \text{for any request $R \in \R^{N \times (M+1)}$}
$$
where
\begin{align*}
e_{i^*} & = \underbrace{(0,0,\dots,0,1,0,\dots,0)}_{\text{$1$ is at position  $i^*$}} \in \R^N \\
\lambda & = (\lambda_0, \lambda_1, \dots, \lambda_M) \in \R^{M+1} \\
i^* & = \argmax_{i=1,2,\dots,N} \lambda_0 R_{i,0} + \lambda_1 R_{i,1} + \lambda_2 R_{i,2} + \dots + \lambda_M R_{i,M}
\end{align*}
Tie breaking in argmax is towards the smaller index $i$.
\end{theorem}
\end{frame}

%%%%%%%%%%%%%%%%%%%%%%%%%%%%%%%%%%%%%%%%%%%%%%%%%%%%%%%%%%%%%%%%%%%%%%%%%%%%%%%%%%%%%%%%%%
\begin{frame}
\frametitle{Lagrangian Policies}

\begin{itemize}
\item Proof of the theorem is too technical.
\item Two basic ingredients:
\begin{itemize}
\item Uniform convergence for Lagrangian policies.
\item Linear programming duality.
\end{itemize}
\item Lagrangian policies are the same as multiclass linear classifiers, or, in other words, multiclass Perceptrons.
\end{itemize}

\end{frame}

\end{document}
